\section{MVP Roadmap}
\label{sec:roadmap}

The roadmap is built to prove one thesis as early as possible: \emph{the same application, unmodified, controls the same class of peripheral on both a legacy machine (via Probe) and a purpose-built rig (via E2B)}. Everything not required to prove that thesis is deliberately deferred, including the FPGA probe front end, the full E2B family, the AP/RT split on real automotive silicon, and any certification-evidence tooling beyond what makes the pipeline auditable in principle.

\subsection{MVP 0 (4--8 weeks)}

\begin{itemize}
\item Single Alfred Bus Module (CAN-FD + LIN) on an STM32/S32K dev board, feeding a laptop-hosted Analyzer.
\item A cheap bench rig with real automotive-style peripherals: a 12V DC motor + Hall sensor (window/actuator stand-in), an LED lamp circuit, a couple of switches, wired to a small breakout that speaks CAN and/or LIN like a toy body-control network (this can literally be a hobbyist CAN-controlled relay/motor board -- the point is a real electrical bus with real physical effect, not a simulator).
\item Passive capture + Level 1/2 decode working end-to-end.
\item Manual Structured Experiment Mode (operator presses switch, script records before/after, computes simple co-occurrence score) --- correlation math implemented, UI can be a notebook/CLI.
\item Demo: Alfred correctly identifies "this CAN ID/bit maps to the lamp" and "this CAN ID maps to the motor," with a confidence score, purely from passive+experiment capture.
\end{itemize}

\subsection{MVP 1 (3 months)}

\begin{itemize}
\item Vehicle Knowledge Graph schema (Section~\ref{sec:knowledge-graph}) implemented as a real, queryable store (not a notebook artifact); Graph-to-IR promotion logic implemented.
\item Machine IR (Section~\ref{sec:machine-ir}) implemented with the Semantic API surface (\code{light.set\_behavior()}, \code{actuator.move\_to\_position()}) callable against the MVP0 rig's promoted Devices.
\item Active Test Mode with the full gating stack from Section~\ref{sec:active-control} (allowlist, state validation, hardware transmit-enable, E-stop, audit log) --- even if the allowlist only ever contains "lamp" and "motor" at this stage, the \emph{mechanism} must be real, not stubbed, because it is the thing that ends up in every later demo unchanged.
\item First E2B-LITE prototype board (real PCB, not a dev-board breakout): MCU + 10BASE-T1S PHY + protected GPIO/PWM/ADC, driving the \emph{same} lamp/switch peripherals from MVP0's rig.
\item Demo: \code{light.set\_behavior(BRAKE)} called against the legacy rig (via Probe/Gateway) and, separately, against the E2B-LITE-controlled version of the same lamp --- same API call, same observable physical result, two different hardware paths underneath.
\end{itemize}

\subsection{MVP 2 (6 months)}

\begin{itemize}
\item E2B-MOTOR prototype board (real PCB): motor driver + current sense + Hall decode, closing a local position-servo loop, exposing \code{LinearActuator} capabilities identical in shape to the legacy-mapped actuator from MVP0/1.
\item Deterministic Firmware Generation pipeline (Section~\ref{sec:firmware-generation}) operating for at least the E2B-LITE and E2B-MOTOR target packs: Machine IR change $\rightarrow$ resource-planned $\rightarrow$ generated firmware $\rightarrow$ built $\rightarrow$ flashed, for a real new Device addition, without hand-written target-specific C code for that addition.
\item Minimal Alfred Lab bench (manual fixtures acceptable): repeatable flash + smoke-test cycle for E2B boards.
\item Engineering file import (DBC at minimum) feeding the Correlation Engine as a confidence prior, demonstrated on a real (or realistic sample/public) DBC against the bench rig.
\item Demo: the full convergence diagram from Section~\ref{sec:convergence} working live --- the same application script drives \code{actuator.set\_position(60\%)} against both the legacy-Probe-mapped actuator and the E2B-MOTOR-controlled actuator, side by side.
\end{itemize}

\subsection{MVP 3 --- first serious customer pilot}

\begin{itemize}
\item Full distributed Probe (multiple synchronized nodes, gPTP) deployed on an actual customer development vehicle (fleet operator, OEM innovation group, or ag/mining equipment maker), Phase 0/1 migration (Section~\ref{sec:migration}) executed for real: passive capture, Vehicle Knowledge Graph built for a meaningful subsystem (e.g., full body/lighting network).
\item Phase 2 (Gateway-mediated control) demonstrated on that customer's actual vehicle for a non-safety-critical function set, under the full Active Test Mode gating stack, with real audit logging reviewed by the customer.
\item Alfred Core first real automotive-silicon build (S32G-class board), running the AP/RT split, as the target for a from-scratch E2B-based subsystem the customer agrees to pilot (e.g., replacing one body ECU) --- this is the first point Phase 3 migration is attempted for real.
\item HIL validation gate operating in the Deployment Manager for that pilot's firmware before anything is flashed to customer hardware.
\item This is the pilot that should convert into a paid engagement and a reference customer, not an internal demo.
\end{itemize}

\begin{decision}[title={Sequencing discipline}]
Do not build Alfred Core (Section~\ref{sec:acc}) or the full E2B family before MVP2/3. Every one of MVP0/1/2's demos is achievable on dev-board-class hardware and a single peripheral class (lamp, then motor); resist the temptation to build the "complete" central compute or probe hardware early, since the thesis-proving demo does not require it and it consumes the runway that MVP0/1's speed depends on.
\end{decision}
